\section{Conceptos fundamentales}

\textbf{Producto de solubilidad ($K_{ps}$):} Es la constante de equilibrio para la disolución de un compuesto iónico poco soluble. Es el producto de las concentraciones de los iones en una solución saturada, cada una elevada a su coeficiente estequiométrico.

\textbf{Efecto del ión común:} La adición de un ión común a una solución saturada disminuye la solubilidad del compuesto, provocando la precipitación del exceso.


\section{Cálculo del producto de solubilidad}

Para un compuesto $A_mB_n$, el equilibrio es:
\[
A_mB_n (s) \rightleftharpoons m A^{n+} (aq) + n B^{m-} (aq)
\]
Y la constante de equilibrio es:
\[
K_{ps} = [A^{n+}]^m [B^{m-}]^n
\]

\resumebox{ejercicio}{Cálculo de $K_{ps}$ y solubilidad}{\footnotesize
    La solubilidad del $AgCl$ en agua es de $1.33 \times 10^{-5}$ M. Calcula su $K_{ps}$. ¿Cuál sería la solubilidad en una solución de NaCl 0.1 M?
}

\resumebox{dato}{Formación de cálculos renales}{\footnotesize
    El producto de solubilidad es importante en el cuerpo humano. La precipitación de sales como el oxalato de calcio ($CaC_2O_4$) en los riñones puede formar cálculos renales. Mantener una adecuada hidratación ayuda a evitar que se alcance el $K_{ps}$ y se produzca la precipitación.
}